\subsection{Sixth letter}

\begin{quotex}
This letter was written shortly after \textbf{Guido de Giorgio}`s visit to \textbf{Rene Guenon} in Blois.

Curiously, the sniping against the young Evola continues. Notice that Guenon treats de Giorgio as a peer while in his letters to Evola, he is always lecturing. Guenon indicates he and de Giorgio are totally in accord regarding their understanding of Tradition. Knowing this, it is clearly a mistake to ignore our translations of de Giorgio and \emph{a fortiori} to ignorantly criticize them as some have.

The Tunisian tariqa of Sheik Ahmad al-Alawi has great significance for Tradition. Besides de Giorgio, Frithjof Schuon was also initiated there. Guenon himself, as this letter shows, was also in communication with that tariqa.

Note, too, that Guenon contributed to the Catholic periodical Regnabit for two years; these articles eventually were included in his Symbols of Sacred Science. Unfortunately, he was forced out by the circle around Jacques Maritain, with whom Guenon had previously engaged in discussions. That has had deleterious consequences for the recovery of Tradition in the West.

\end{quotex}
Paris, 1 November 1927

We had learned with pleasure, from your card sent on your arrival at Varazze, that your return trip went quite well; but following your letter, which reached us before our departure from Blois, brought us unfortunately much less good news. However, as you don't speak again of your health in your last letter, we like to believe that it has now been restored.

In spite of everything that I know of Evola, especially from you, I was a little surprised by his refusal to use your article\footnote{For Ur.}. I ask myself, under these conditions, why he is so insistent that I send him something, because he must certainly think what I would write would also be totally traditional, and consequently, would not satisfy him at all. I see that obviously there is nothing to do with him; also, I am going to write him very gladly in the way that you ask me, all the more reason so that will cut short (at least I think so) a new insistence of his part and will give me a decisive reason to not collaborate in his journal. Besides, I owe him a response, having received from him, about a month ago, a letter in which he asked me for some information on a point which I believe that he had already spoken to you about. I am going to enclose my letter with this one, while asking you to then send it to him from there. I apologize for being unable to stamp it, not having any Italian stamps at my disposal.

I would be happy to read your article if Evola returned it to you and if you want to send it to me; but I am sure in advance that I would be completely in accord with you.

What is the Latin text for which Reghini's translation seems doubtful to you?

I am curious to know what Evola might say about Milarepa; he will certainly be scandalized by the totally inferior place that is given to magic, however well it suits him in reality.

I do not know at all Ashtavakra's poem\footnote{Ashtavakra Gita} that you mention. If you would like, as you tell me, to send me the translation that you have, although it is rather an adaptation, that will interest me and will give me at least an idea of what the poem is about.

Here is the address you asked about: \footnote{G. provides the Tunisian address of Eugene Taillard, who was affiliated with the tariqa of Sheik Ahmad al-Alawi.}

The booklet that [Gustave-Henri] Jossot sent me hardly contains anything of interest except the summary of some interviews with Kheireddine and an outline on the doctrine of the Alouites to which he is now associated as I told you. It seems that this brotherhood has been expanding; I just learned that it has a zawiya\footnote{Islamic religious school.} in Paris, on Blvd St Germain, close to here. That makes me fear that it may become too open and deviate just like so many others. It also appears that it is working, in Algeria, to promote a rapprochement between Christians and Muslims, which is definitely very laudable, but will these efforts lead to an appreciable result?

The journal Parise alluded to is Regnabit. I actually collaborated in it regularly from August 1925 until May 1927, and I had always forgotten to tell you about it. But see how history is told: this journal is as little directed by the Jesuits as they did everything they could to prevent it from appearing. Moreover, I had to end this collaboration following certain machinations which had their origins in [Jacques] Maritain's entourage. It would take too long to tell you about it and moreover rather less interesting au fond. I intend to take up again someday the questions of symbolism that I had begun to treat in that journal in order to create a book out of it; but when will I get to it?

I finished correcting the proofs of the Crisis of the Modern World. I think that it will appear around mid-November.



\flrightit{Posted on 2012-09-17 by Cologero }

\begin{center}* * *\end{center}

\begin{footnotesize}\begin{sffamily}



\texttt{Graham on 2012-09-17 at 21:49 said: }

" … without forgetting the little piece of perfectly gratuitous sarcasm …"

His opinion of Evola has been consistent low over the course of this correspondance. But this was during the 20s. Do we know if it changed during the 30s and 40s, and if so, for what reason, or reasons?


\hfill

\texttt{Cologero on 2012-09-17 at 22:23 said: }

We have documented much of this already. Guenon gave a generally positive review of ``Revolt against the modern World" after its publication. After the 1920s, there was limited contact between Evola and Guenon aside from exchanging books or articles. Evola translated Guenon's works in the 30s. Correspondence resumed after Evola's recuperation; we have made several of those letters available. Again, there was some disagreement between them, but relations were cordial.


\end{sffamily}\end{footnotesize}
